% Source-derived chapter 16: Proof of the main result
% Original source: no-enriques-surface-over-the-integers
\chapter{Proof of the main result}
\label{no-enriques-surface-over-the-integers-chapter-16}
\source{no-enriques-surface-over-the-integers}

\begin{remark}[Source section]
\label{no-enriques-surface-over-the-integers-chapter-16-source}
\source{no-enriques-surface-over-the-integers}
This chapter records the source section; no Lean declaration has yet been checked for this material.
\notready
\end{remark}

% The source section title was: Proof of the main result
\mylabel{Proof}

We are finally ready to prove Theorem \ref{no enriques over integers},  the main result of this paper.
The task is to show that the fiber category $\shM_\Enr(\ZZ)$ for the stack of Enriques surfaces is empty.
Seeking a contradiction, we assume that that there is a family of Enriques surfaces
$f:\foY\ra \Spec(\ZZ)$. According to Proposition \ref{constant picard scheme}, the Picard scheme $\Pic_{\foY/\ZZ}$ is constant.
So the closed fiber $Y=\foY\otimes\FF_2$ also has constant Picard scheme.
It is  non-exceptional, by  \cite{Schroeer-2021b}, Theorem 7.2, because this Enriques surface
 comes with a deformation over $W_2(\FF_2)=\ZZ/4\ZZ$.
Therefore it suffices to establish:

\begin{theorem}
\source{no-enriques-surface-over-the-integers}
\notready
\mylabel{no non-exceptional}
There is no  Enriques surface $Y$ over the prime field $k=\FF_2$  that is  non-exceptional and  whose   Picard scheme
is constant.
\end{theorem}

\proof
Suppose such a $Y$ exists.
By Theorem \ref{family with constant pic} there is a genus-one fibration $\varphi:Y\ra\PP^1$.
For each geometric point $\ba:\Spec(\Omega)\ra \PP^1$ lying over a non-rational closed point  $a\in\PP^1$, the geometric fiber $Y_\ba$
is  irreducible, according to Proposition \ref{constant num over finite field},  thus the  Kodaira symbol is $\I_0, \I_1$ or $\II$.
Moreover, it must be simple by Theorem \ref{family with constant pic}.
The possible configurations of fibers over the rational points $t=0,1,\infty$ are tabulated
in Proposition \ref{key reduction}.
 In Sections \ref{Elimination I4*} and \ref{Elimination III*} we saw that the Kodaira symbols
$\I_4^*$ and $\III^*$ actually do not occur.
For the classical Enriques surface $\bY=Y\otimes_kk^\alg$   this means that for every genus-one fibration, the configuration
of degenerate fibers is either $\I_1^*+\I_4$ or $\I_2^*+\III+\III$. Furthermore, in the former case the $\I_1^*$-fiber is multiple,
whereas in the latter case two of the three fibers are multiple.
 We showed in Proposition \ref{only two configurations} that
such an  Enriques surface does not exist. This gives the desired contradiction.
\qed


%===========================================================
